\documentclass[12pt]{article}
\usepackage{graphicx} % Required for inserting images
\usepackage[margin=0.5in]{geometry}
\usepackage{amsmath, amssymb}
\usepackage{mathrsfs}


\usepackage{soul}


% Dr. David Brown Commands
\newcommand{\ds}{\displaystyle}
\newcommand{\R}{\mathbb{R}}
\newcommand{\N}{\mathbb{N}}
\newcommand{\Q}{\mathbb{Q}}
\newcommand{\C}{\mathbb{C}}


% Commands to get script r
\usepackage{calligra}
\DeclareMathAlphabet{\mathcalligra}{T1}{calligra}{m}{n}
\DeclareFontShape{T1}{calligra}{m}{n}{<->s*[2.2]callig15}{}
\newcommand{\scripty}[1]{\ensuremath{\mathcalligra{#1}}}

% Unit commands
\newcommand{\degree}{\text{°}}
\newcommand{\unitSpeed}{\frac{\text{m}}{\text{s}}}
\newcommand{\unitAcceleration}{\frac{\text{m}}{\text{s}^2}}

% Constant commands
\newcommand{\avogadro}{6.022\times10^{23}}

\newcommand{\boltzmannConstK}{1.381\times10^{-23}\frac{\text{J}}{\text{K}}}
\newcommand{\boltzmannConstInEV}{8.617\times10^{-5}\frac{\text{eV}}{\text{K}}}
\newcommand{\planckConst}{6.626\times10^{-34}\text{J}\cdot\text{s}}

\newcommand{\atmP}{1.01325\times10^5\,\text{Pa}}

% Commands to easily write partial derivatives
\newcommand{\partDeriv}[2]{\frac{\partial #1}{\partial #2}}
\newcommand{\doublePartDeriv}[2]{\frac{\partial^2 #1}{\partial^2 #2}}


\title{Problem 7.7}
\author{Gabriel Decker}


\begin{document}



\maketitle
\begin{flushleft}



In this problem, we will consider various situations involving distinguishable and indistinguishable particles in a Boltzmann-statistics system.
These particles can have one of 10 single-particle states which all have energy zero.
This implies that the Boltzmann factor for each state is $e^{-E(s)/kT}=e^0=1$.
Therefore, the partition function is simply the number of available single-particle states.\\~\\

\textbf{Part a:}\\
As stated above, the partition function $Z$ is simply the number of possible states.
Since we have one particle with 10 states available to it, $Z=10$.\\~\\

\textbf{Part b:}\\
Since there are now two distinguishable particles, and each particle has 10 possible states, this implies that $Z=10\cdot10=100$.\\~\\

\textbf{Part c:}\\
We now consider two indistinguishable bosons.
For this, there are two possible general cases for states.
The particles can either be in different states or the same states.
Since there are 10 states, there are 10 possible ways for them to be in the same state.
Now, consider the case where the particles are in different states.
Choose the state for each particle sequentially.
For the first particle, there are 10 states to choose from.
For the second particle, there are 9 states since we're avoiding the state the first particle is in.
Therefore, there are $10\cdot9=90$ states for this case.
However, since our particles are indistinguishable, we must divide this by 2.
This is because we would otherwise double-count these states.
So, there are $90/2=45$ states in this case.
This brings the partition function to $Z=45+10=55$.\\~\\

\textbf{Part d:}\\
We now consider two indistinguishable fermions.
This is the same as for the bosons, except states where both particles share a single-particle state are eliminated.
So, $Z=45$.\\~\\

\textbf{Part e:}\\
The relevant equation is the following.
\begin{equation}
    Z=\frac{1}{N!}Z_1^N
\end{equation}
This equation is only true for indistinguishable particles that are only possibly in different states. It is simply an approximation for not-very-dense systems ($Z_1\gg N$).
For our situation, $Z_1$ is simply what we computed in part a and $N=2$.
So, we get the following.
$$Z=\frac{1}{N!}Z_1^N$$
$$\implies Z=\frac{1}{2!}10^2$$
$$\implies Z=\frac{1}{2}100$$
$$\implies Z=50$$
This is between the fermion and boson case.\\~\\

\textbf{Part f:}\\
Now, we will use the partition functions (or the number of possible states, in this case) for to calculate the probabilities of both particles being in the same state.
This will be done for the three scenarios: distinguishable particles, bosons, and fermions.\\~\\

For distinguishable particles, the total number of states is $Z=100$.
There are 10 states where the particles are in the same state.
Therefore, $\mathcal{P}=10/100=1/10=0.1$.\\~\\

For bosons, the total number of states is $Z=55$.
There are 10 states where the particles are in the same state.
Therefore, $\mathcal{P}=10/55=0.182$.\\~\\

For fermions, the total number of states is $Z=45$.
However, it is impossible for fermions to share a state, so $\mathcal{P}=0/45=0$.\\~\\





\end{flushleft}

\end{document}
